\documentclass[11pt]{article}
\usepackage[letterpaper,margin=0.95in]{geometry}
\usepackage{amsmath,amssymb,amsthm,mathtools,booktabs,microtype}
\usepackage[T1]{fontenc}
\usepackage{lmodern}
\usepackage[colorlinks=true,linkcolor=blue!45!black,citecolor=blue!45!black,urlcolor=blue!55!black]{hyperref}
\usepackage{xcolor}
\usepackage{fancyhdr}
\hypersetup{pdftitle={A constant-free binomial filter for prime discrepancies},pdfauthor={Research note}}
\pagestyle{fancy}\fancyhf{}\lhead{\small\textsc{Prime discrepancies and Dirichlet zeros}}\rhead{\small November 2026}\cfoot{\thepage}\setlength{\headheight}{14pt}
\newtheorem{theorem}{Theorem}[section]
\newtheorem{lemma}[theorem]{Lemma}
\newtheorem{corollary}[theorem]{Corollary}
\newtheorem{proposition}[theorem]{Proposition}
\theoremstyle{remark}\newtheorem{remark}[theorem]{Remark}
\newcommand{\R}{\mathbb R}
\newcommand{\C}{\mathbb C}
\newcommand{\N}{\mathbb N}
\newcommand{\Z}{\mathbb Z}
\newcommand{\ord}{\operatorname{ord}}
\newcommand{\Rea}{\operatorname{Re}}
\newcommand{\dist}{\operatorname{dist}}
\newcommand{\dd}{\,\mathrm d}
\newcommand{\GRH}{\mathrm{GRH}}
\title{\vspace{-1em}\bfseries A constant-free binomial filter\\ for prime discrepancies\\[5pt]\large Uniform central-zero recovery and a robust one-sided GRH criterion}
\author{Research note}
\date{26 November 2026}
\begin{document}
\maketitle
\begin{abstract}
We construct a finite, positive binomial average of a difference of ramp-weighted von Mangoldt sums. For every fixed modulus and every subgroup of its unit group, assuming GRH for the associated Dirichlet characters, this average converges uniformly over all starting scales to four times the total multiplicity of their central zeros, at rate $O(\log(2m)/\sqrt m)$ after $m$ repetitions. Neither the centering constants nor any zero ordinates are required by the statistic. Conversely, for any fixed filter depth, an eventual upper bound for the filtered discrepancy excludes all off-line zeros of the associated family. The proof combines the smoothed prime remainder and one-sided Landau principle emphasized in the recent Lean GRH development with a two-scale, badly approximable translation filter. The result also isolates central multiplicities of individual real characters by selecting a subgroup of residue classes. This is an independently derived candidate contribution, not a claim of established priority or a new machine-checked proof of GRH.
\end{abstract}

\section{Input, motivation and scope}
The repository \texttt{lsernau/grh-lean} \cite{repo} is a shortened rewriting of the Alp\"oge--Claude Lean development. We inspected its README, reading guide, statement file, main theorem file and relevant compact endgame at commit
\begin{center}\small\texttt{72fead102f663d5c7a407964a13dc0cf9d977207}.\end{center}
The five top-level statements include Mathlib's \texttt{RiemannHypothesis}, nonvanishing to the right of $1/2$ (excluding the principal pole), and the critical-line statement for all Dirichlet characters, primitive or not. The README reports standard-axiom dependencies and successful machine checks; independent checks of the underlying development are documented by Rubinstein \cite{rubinstein}. No full local Lean replay is performed here. We prove the implication from the stated Dirichlet GRH; it is unconditional relative to acceptance of that input. No GRH for arbitrary number fields or automorphic forms is used.

In \texttt{grh/src/New/Reduction/OneSidedBound.lean}, the key quantity is a centered sum with weight $(1/n-1/x)$ over $n\le x$, restricted to $n\equiv1\pmod q$ and multiplied by a quadratic nonsplitting weight. The compact \texttt{src/New/Main.lean} combines bounds on integral-form volumes and indices in \texttt{star\_inequality}, yielding an eventual one-sided prime-error bound. The analytic endgame in \texttt{New/Endgame/OneSidedCriterion.lean} then applies a Landau argument. The Wikipedia overview \cite{wiki} highlights the relation between zeros, the prime-counting error and Dirichlet progressions. Its broader uses of GRH must be distinguished from the Dirichlet theorem actually provided by the repository.

Our construction does not attempt to transfer the high-rank theta-lift machinery to another family. Instead, it transfers the analytic ramp and one-sided argument, and adds a discrete filtering principle. We ask whether the oscillatory contributions can be strongly attenuated while retaining both the zero-free-half-plane criterion and the central-zero signal. The two objectives would conflict for a generic averaging filter; the special multiplier below resolves this conflict.

\paragraph{Related work and novelty.} One-sided Mellin criteria for RH have a long history (Landau; Montgomery--Vaughan \cite{mv}). Suzuki \cite{suzuki} treats sign conditions and weighted summatory functions, including the role of central zeros; Platt--Ramar\'e \cite{pr} study the centered $\Lambda(n)/n$ sum in arithmetic progressions. Explicit formulae and zero counts are classical \cite{dav,bennett}. Our proposed contribution is the conjunction of a constant-eliminating difference, a fixed two-step positive binomial average independent of the modulus and zeros, a uniform quantitative central-multiplicity limit, and preservation of a one-sided criterion at every fixed depth. Searches for binomial smoothing of prime discrepancies, logarithmic averaging of central zeros and one-sided weighted prime criteria located the works above, but did not establish that this precise construction has appeared before. A thorough literature review is still needed; priority is not asserted.

\section{Statement of results}
Let $q\ge1$, $G=(\Z/q\Z)^\times$ (trivial when $q=1$), and $H\le G$. Write
\[
 X=H^\perp=\{\chi\bmod q:\chi(h)=1\ \text{for every }h\in H\}.
\]
All characters are extended by zero on nonunits; for $q=1$ the character is identically one. Define
\begin{align}
 a(n)&=[G:H]\,\Lambda(n)\,\mathbf 1_{n\bmod q\in H},\nonumber\\
 S(x)&=\sum_{n\le x}a(n)\left(\frac1n-\frac1x\right),\qquad x\ge1,\label{eq:S}\\
 d(t)&=\frac{e^{(t+1)/2}}{e^{1/2}-1}\{1+S(e^t)-S(e^{t+1})\},\qquad t\ge0.\label{eq:d}
\end{align}
For a function $f$ on $[0,\infty)$, let $T_hf(t)=f(t+h)$ and $A_h=(I+T_h)/2$. Our filter of depth $m$ is
\begin{equation}\label{eq:B}
 B_m=A_1^m A_{\sqrt2}^m,\qquad
 B_m f(t)=4^{-m}\sum_{j,k=0}^m\binom mj\binom mk f(t+j+k\sqrt2).
\end{equation}
The $(m+1)^2$ distinct shifts have nonnegative weights summing to one. Set
\[
 M=\sum_{\chi\in X}\ord_{s=1/2} L(s,\chi),
\]
where the order is zero for a nonvanishing value. Let $\GRH_X$ mean that every zero of these functions in $0<\Rea s<1$ lies on $\Rea s=1/2$.

\begin{theorem}[Uniform central-multiplicity recovery]\label{th:main}
Assume $\GRH_X$. There is a constant $C_{q,H}$ such that for all integers $m\ge1$ and all real $t\ge0$,
\begin{equation}\label{eq:main}
 |B_m d(t)-4M|\le C_{q,H}\frac{\log(2m)}{\sqrt m}.
\end{equation}
In particular, $B_m d(t)\to4M$ uniformly in $t\ge0$. The statistic in \eqref{eq:d}--\eqref{eq:B} uses only finite prime-power sums, elementary constants, and no $L$-value, centering constant, or zero ordinate. For sufficiently large $m$ (depending on the family), $M$ is the nearest integer to $B_m d(t)/4$, simultaneously for all $t\ge0$.
\end{theorem}

\begin{theorem}[One-sided detection, without GRH]\label{th:converse}
For every fixed integer $m\ge0$, the following statements are equivalent:
\begin{enumerate}
\item $\GRH_X$.
\item $B_m d(t)$ is bounded above for all sufficiently large $t$.
\item $B_m d(t)$ is bounded in absolute value for all sufficiently large $t$.
\end{enumerate}
More quantitatively, if an associated $L$-function has a zero $\rho$ with $\Rea\rho>1/2+\delta$, where $\delta\ge0$, then for every fixed $m\ge0$,
\begin{equation}\label{eq:blowup}
 \limsup_{t\to\infty}e^{-\delta t}B_m d(t)=+\infty.
\end{equation}
\end{theorem}

Taking $H=\{1\}$ makes $X$ the full group of characters modulo $q$ and $[G:H]=\varphi(q)$. Thus the ordinary identity residue class alone suffices for the criterion. Taking $H=G$ gives the principal character alone. All constants in the results are for a fixed family, with no uniformity in $q$ claimed.

\section{Analytic preparation: the ramp and its poles}
Character orthogonality, including the extension by zero, gives $a(n)=\Lambda(n)\sum_{\chi\in X}\chi(n)$. Put
\begin{equation}\label{eq:F}
 F(s)=\sum_{n\ge1}\frac{a(n)}{n^s}=-\sum_{\chi\in X}\frac{L'}{L}(s,\chi),\qquad \Rea s>1.
\end{equation}
The right side meromorphically continues. Exactly one principal character lies in $X$. By the classical nonvanishing on $\Rea s=1$ away from its pole, the Laurent expansion at $1$ is
\[
 F(1+z)=\frac1z+c_0+O(z).
\]
The constant $c_0$ is real by conjugation symmetry. Let $c=c_0-1$ and set
\begin{equation}\label{eq:R}
 R(x)=S(x)-\log x-c,\qquad r(t)=e^{t/2}R(e^t).
\end{equation}
These definitions are unconditional. No estimate for $R$ is needed to derive its transform. The bound $a(n)\le [G:H]\log n$ gives $S(e^t)=O_{q,H}((1+t)^2)$, so $r(t)=O_{q,H}(e^{t/2}(1+t)^2)$.

\begin{lemma}[Laplace transform]\label{lem:laplace}
The Laplace transform $\mathcal Lr(w)$ initially for $\Rea w>1/2$ has the meromorphic continuation
\begin{equation}\label{eq:laplace}
 \mathcal Lr(w)=\frac{F(w+1/2)}{(w-1/2)(w+1/2)}
   -\frac1{(w-1/2)^2}-\frac c{w-1/2}.
\end{equation}
The singularity at $w=1/2$ is removable. If $\rho$ is a zero in the critical strip of at least one member of $X$, and $b(\rho)=\sum_{\chi\in X}\ord_\rho L(s,\chi)$, the residue at $w=\rho-1/2$ is
\begin{equation}\label{eq:residue}
 \frac{b(\rho)}{\rho(1-\rho)}.
\end{equation}
In particular a real zero $1/2<\rho<1$ contributes a strictly positive residue. There are no other singularities in $\Rea w>0$.
\end{lemma}
\begin{proof}
For $\Rea z>0$, absolute convergence permits interchange of the sum and the integral:
\begin{align*}
 \int_0^\infty S(e^t)e^{-zt}\dd t
 &=\sum_{n\ge1}a(n)\int_{\log n}^\infty (n^{-1}-e^{-t})e^{-zt}\dd t\\
 &=\frac{F(1+z)}{z(z+1)}.
\end{align*}
Subtract $1/z^2+c/z$ and substitute $z=w-1/2$. The expansion $F(1+z)/(z(z+1))=z^{-2}+(c_0-1)z^{-1}+O(1)$ proves removability. At a zero the residue of $F$ is $-b(\rho)$; all character weights are $+1$, so coincident zeros add, never cancel. The denominator gives \eqref{eq:residue}. The Euler product and classical nonvanishing on the line $1$ exclude further singularities in the stated half-plane. Imprimitive Euler factors have their zeros on $\Rea s=0$ only and are irrelevant here.
\end{proof}

The unknown $c$ disappears exactly via the translation operator
\begin{equation}\label{eq:D}
 d(t)=\frac{e^{1/2}I-T_1}{e^{1/2}-1}\,r(t).
\end{equation}
Indeed substitution of \eqref{eq:R} yields \eqref{eq:d}. This identity is useful because the translation operator preserves constants while annihilating $ce^{t/2}$, the nuisance mode associated with centering. An equivalent prime-distribution interpretation is obtained by setting $\Psi(u)=\sum_{n\le u}a(n)$. The ramp is continuous and piecewise differentiable, with derivative $S'(u)=\Psi(u)/u^2$ away from the jump points of $\Psi$. Integration therefore gives
\[
 d(\log x)=\frac{\sqrt{ex}}{\sqrt e-1}
 \int_x^{ex}\frac{u-\Psi(u)}{u^2}\dd u.
\]
Thus the constant-free statistic is a normalized integral of the deficit in the weighted Chebyshev count over a multiplicative interval.

\begin{lemma}[Absolutely convergent spectral expansion]\label{lem:spectral}
Assume $\GRH_X$. Let $\Gamma$ be the multiset of ordinates of all nontrivial zeros of the members of $X$ (with multiplicity, including both signs and any zeros at $0$). Then, uniformly for $t\ge0$,
\begin{equation}\label{eq:explicit}
 r(t)=\sum_{\gamma\in\Gamma}\frac{e^{i\gamma t}}{1/4+\gamma^2}+E(t),\qquad
 |E(t)|\le C_{q,H}e^{-t/4}.
\end{equation}
The series is absolutely and uniformly convergent. The count of ordinates of absolute value at most $T$ is $O_{q,H}(T\log(2+T))$ for $T\ge1$.
\end{lemma}
\begin{proof}
We give contour details so that no unsmoothed explicit-formula convergence convention enters. The standard Dirichlet $L$-function zero-counting formula (with multiplicity) gives $N(T)=O_{q,H}(T\log(2+T))$ and $N(T+1)-N(T-1)=O_{q,H}(\log(2+T))$ for $T\ge2$; see \cite{dav,bennett}. The conductor-one case uses the classical zeta zero count. Imprimitive factors add no zeros inside the critical strip. Dyadic summation gives $\sum_{\gamma\in\Gamma}(1/4+\gamma^2)^{-1}<\infty$.

The elementary inverse Mellin identity for the ramp is
\[
 \frac1{2\pi i}\int_{c_1-i\infty}^{c_1+i\infty}\frac{y^z}{z(z+1)}\dd z
 =\begin{cases}1-y^{-1}&y\ge1,\\0&0<y<1,\end{cases}\quad c_1>0.
\]
It follows by closing the contour to the left or right; at $y=1$ the continuous ramp has value zero. Equivalently, it is Mellin inversion of the function on the right; the transform is integrable on the given line. Apply it termwise at $c_1=1$ using absolute convergence of \eqref{eq:F}. For $x\ge1$,
\begin{equation}\label{eq:perron}
 S(x)=\frac1{2\pi i}\int_{1-i\infty}^{1+i\infty}
       \frac{F(1+z)x^z}{z(z+1)}\dd z.
\end{equation}
Shift to $\Rea z=-3/4$ (equivalently $\Rea s=1/4$). Here are the bounds justifying the shift. For $|v|\ge2$ and $1/4\le\sigma\le2$, the usual local partial-fraction estimate is
\[
 F(\sigma+iv)=-\sum_{\substack{\rho:\,|\Im\rho-v|\le1}}\frac{b(\rho)}{\sigma+iv-\rho}
       +O_{q,H}(\log(2+|v|)),
\]
where the sum is over distinct nontrivial zeros. One can obtain this estimate directly by subtracting the completed Hadamard logarithmic derivative at $2+iv$ from that at $\sigma+iv$. The far-zero differences are $O((1+|v-\Im\rho|)^{-2})$, summing to $O_{q,H}(\log(2+|v|))$ by the unit-interval zero count. For the near zeros the terms at $2+iv$ are bounded, the gamma logarithmic derivatives contribute at most $O_{q,H}(\log(2+|v|))$, and the Euler-product logarithmic derivative at $2+iv$ is bounded. The finitely many imprimitive Euler factors have bounded logarithmic derivatives for $\sigma\ge1/4$. These arguments also apply to the completed zeta function after separating its elementary pole factors.

By GRH, the distance from $\sigma=1/4$ to every nontrivial zero is at least $1/4$. The vertical line therefore has $F(1/4+iv)=O_{q,H}(\log(2+|v|))$ (bounded on compact parts). For arbitrarily large $T\to\infty$, both heights $\pm T$ can be chosen at distance $\gg_{q,H}1/\log(2+T)$ from all ordinates: exclude intervals of that radius around the $O_{q,H}(\log T)$ ordinates near each of $\pm[T,T+1]$, choosing the constant small. On these horizontal segments $F=O_{q,H}(\log^2(2+T))$. Their integrals in \eqref{eq:perron} tend to zero for each fixed $x$, since the denominator has size $\asymp T^2$.

The crossed poles are $z=0$ with residue $\log x+c$, and $z=\rho-1$ with residue $b(\rho)x^{\rho-1}/[\rho(1-\rho)]$. The new line contributes $O_{q,H}(x^{-3/4})$, with a constant independent of $x\ge1$, by absolute integrability of its logarithmic bound divided by $|z(z+1)|$. Summing residues is legitimate by the absolute convergence already shown. Multiplication by $x^{1/2}$ and $x=e^t$ gives \eqref{eq:explicit}.
\end{proof}

\section{The positive translation filter}
For an exponential mode $e^{wt}$, the multiplier of $A_h$ is $(1+e^{hw})/2$. The multiplier of \eqref{eq:D} is
\begin{equation}\label{eq:K}
 K(w)=\frac{e^{1/2}-e^w}{e^{1/2}-1}.
\end{equation}
Thus the combined multiplier is
\begin{equation}\label{eq:mult}
 P_m(w)=K(w)\left(\frac{1+e^w}{2}\frac{1+e^{\sqrt2 w}}{2}\right)^m.
\end{equation}
It has three important properties:
\begin{itemize}
\item $P_m(0)=1$, preserving the central term.
\item $P_m(w)\ne0$ whenever $0<\Rea w<1/2$. The Bernoulli factors have zeros only on $\Rea w=0$; the zeros of $K$ are on $\Rea w=1/2$.
\item $P_m(u)>0$ for real $0<u<1/2$, preserving the sign of real off-line residues.
\end{itemize}
On the imaginary axis the Bernoulli part has modulus
\begin{equation}\label{eq:cos}
 \big|\cos(\gamma/2)\cos(\sqrt2\gamma/2)\big|^m.
\end{equation}
Neither $K(i\gamma)$ nor the original denominator impedes convergence: $|K(i\gamma)|\le(e^{1/2}+1)/(e^{1/2}-1)$, and the denominator is quadratic in $\gamma$.

\begin{lemma}[Two-scale damping]\label{lem:cos}
For every $\eta>0$ there is $c_\eta>0$ such that for all real $|u|\ge\eta$,
\begin{equation}\label{eq:damping}
 |\cos(u/2)\cos(\sqrt2 u/2)|\le
 \exp\left(-\frac{c_\eta}{(1+|u|)^2}\right).
\end{equation}
\end{lemma}
\begin{proof}
Write $v=u/(2\pi)$, choose nearest integers $a$ to $v$ and $b$ to $\sqrt2v$, and put $d_1=|a-v|$, $d_2=|b-\sqrt2v|$. For $|u|$ sufficiently large, $(a,b)\ne(0,0)$, so the nonzero integer $b^2-2a^2$ has absolute value at least one. Hence
\[
 1\le|b^2-2a^2|\le (d_2+\sqrt2d_1)(|b|+\sqrt2|a|)
       \ll (d_1+d_2)(1+|u|).
\]
It follows that $d_1^2+d_2^2\gg(1+|u|)^{-2}$. For the bounded range $\eta\le|u|\le U$, the continuous sum of squared distances has strictly positive minimum: simultaneous integer values of $v,\sqrt2v$ force $v=0$. Finally, for $0\le d\le1/2$,
\[
 |\cos(\pi d)|=1-2\sin^2(\pi d/2)\le1-2d^2\le e^{-2d^2}.
\]
The sine bound follows by concavity on $[0,\pi/2]$. Multiplying the two bounds proves \eqref{eq:damping}.
\end{proof}

\begin{lemma}[Summing the damping]\label{lem:sum}
For a discrete multiset $\Gamma$ of real ordinates with finite count on compact intervals and $N(T)=O(T\log(2+T))$, and with any fixed $c_*>0$, one has
\[
 \sum_{\substack{\gamma\in\Gamma\\\gamma\ne0}}
 \frac{\exp(-c_*m/(1+|\gamma|)^2)}{1/4+\gamma^2}
 \ll_{\Gamma,c_*}\frac{\log(2m)}{\sqrt m},\qquad m\ge1.
\]
\end{lemma}
\begin{proof}
The terms $0<|\gamma|<1$ are finite and contribute $O(e^{-c_*m/4})$. Set $Y=\sqrt m$. For $|\gamma|\ge Y$, split into intervals $[2^jY,2^{j+1}Y)$, $j\ge0$. The sum without the exponential is $\ll\sum_{j\ge0}(j+\log(2Y))/(2^jY)\ll\log(2Y)/Y$. For $1\le|\gamma|<Y$ use intervals $[Y/2^{j+1},Y/2^j)$, intersected with $[1,Y)$. In each nonempty interval the upper endpoint $U=Y/2^j$ is at least $1$, the denominator is $\gg U^2$, and $1+|\gamma|\le2U$. The interval contributes at most
\[
 C\frac{\log(2Y)}Y\,2^j e^{-c_*4^j/4}.
\]
This is summable over $j\ge0$. Combining the three ranges gives the result.
\end{proof}

\begin{proof}[Proof of Theorem \ref{th:main}]
The multiset of nonzero ordinates has a positive gap from zero (or is empty), by isolated zeros and finiteness on compact sets for a finite family. Choose $\eta$ smaller than that gap. Apply \eqref{eq:D} and \eqref{eq:B} termwise to \eqref{eq:explicit}; this is valid for each $m$ by absolute uniform convergence. The zero ordinate occurs $M$ times and contributes $4M$ since $P_m(0)=1$. For other ordinates use \eqref{eq:cos}, Lemma \ref{lem:cos} and Lemma \ref{lem:sum} to get a uniform $O_{q,H}(\log(2m)/\sqrt m)$ bound.

The error in \eqref{eq:explicit} after \eqref{eq:D} remains $O_{q,H}(e^{-t/4})$. Each positive average $A_h$ replaces that envelope by $(1+e^{-h/4})/2$ times the envelope. Thus the final error is bounded by $C_{q,H} e^{-t/4}\kappa^m$, where
\[
 \kappa=\frac{1+e^{-1/4}}2\frac{1+e^{-\sqrt2/4}}2<1.
\]
This is uniform also at $t=0$. The two estimates prove \eqref{eq:main}. Eventually the right side is strictly less than $2$, which makes rounding after division by $4$ unambiguous. This is an existence statement, not a claimed explicit numerical threshold in $q,H$.
\end{proof}

\section{Why hypothetical off-line zeros cannot disappear}
For $h\ge0$ and $r$ as above, the Laplace transform satisfies
\[
 \mathcal L(T_hr)(w)=e^{hw}\left(\mathcal Lr(w)-\int_0^h r(v)e^{-wv}\dd v\right).
\]
The finite-interval integral is entire in $w$. Consequently the transform of $B_md$ equals $P_m(w)\mathcal Lr(w)$ plus an entire function. Every off-line zero in $\Rea s>1/2$ gives a genuine pole after filtering by \eqref{eq:mult}; a real such zero gives a positive residue. This argument does not assume GRH, zero simplicity, or distinct ordinates between characters.

\begin{lemma}[Landau principle for Laplace integrals]\label{lem:landau}
Let $f$ be locally integrable and nonnegative on $[T,\infty)$ and assume its Laplace integral converges for sufficiently large real $w$. If its abscissa of convergence $\alpha$ is finite, then the integral cannot extend holomorphically across the real point $w=\alpha$.
\end{lemma}
\begin{proof}
Suppose it did. Pick a small disk about $\alpha$ of radius $r>0$ on which the extension is holomorphic, and a real $s_0=\alpha+r/4$. The Taylor series in $z$ for the integral evaluated at $s_0-z$ converges for $|z|<3r/4$. We may assume $T\ge0$ by discarding a finite interval, whose integral is entire. Differentiation under the integral for $s_0>\alpha$ gives Taylor coefficients
\[
 \frac1{k!}\int_T^\infty t^k f(t)e^{-s_0t}\dd t\ge0.
\]
At $z=r/2$, monotone convergence shows the series equals the integral at $s_0-r/2<\alpha$ and is finite. This contradicts the definition of the abscissa.
\end{proof}

\begin{proof}[Proof of Theorem \ref{th:converse}]
Under $\GRH_X$, Lemma \ref{lem:spectral} and the boundedness of the finite operators show boundedness of $B_m d$ for every $m\ge0$, even before taking $m\to\infty$.

Conversely, let $g=B_m d$ with $m$ fixed, and suppose $g(t)\le C e^{\delta t}$ eventually for some $\delta\ge0$ and real $C$. If $\rho$ is a zero with $\Rea\rho>1/2+\delta$, the transform of $g$ has a pole at $w_0=\rho-1/2$. Define $f(t)=Ce^{\delta t}-g(t)$ on a tail $t\ge T$ where it is nonnegative. Its Laplace transform is, initially on a far-right half-plane, $C/(w-\delta)-P_m(w)\mathcal Lr(w)$ plus an entire function (absorbing the finite initial interval). It has a pole at $w_0$, so its abscissa satisfies $\alpha\ge\Rea w_0>\delta\ge0$. Otherwise its holomorphic integral on $\Rea w>\alpha$ would contradict that pole by analytic continuation. The growth estimate for $r$, and the fact that the shifts are fixed, give $\alpha\le\max(1/2,\delta)=1/2$ in the case under consideration; $\delta<1/2$ follows from the zero-free line $1$.

By Lemma \ref{lem:landau}, the transform must be singular at the real $w=\alpha$. The expression is meromorphic there. If it is regular (in particular if $\alpha=1/2$), we have a contradiction. Otherwise $0<\alpha<1/2$ and the singularity is a simple pole arising from a real zero $\rho'=\alpha+1/2$. Its residue for $\mathcal Lf$ is
\[
 -P_m(\alpha)\frac{b(\rho')}{\rho'(1-\rho')}<0.
\]
Approaching $\alpha$ from above on the real axis, the transform tends to $-\infty$, impossible for a Laplace integral of a nonnegative function. No eventual upper bound is possible. This proves \eqref{eq:blowup}, and the case $\delta=0$ shows that upper boundedness excludes every zero to the right of the line. The primitive functional equation reflects any left-hand zero into a right-hand zero of a conjugate character, still in $X$; induction to modulus $q$ changes no zeros in the open strip. Hence $\GRH_X$ follows. Absolute boundedness implies upper boundedness trivially.
\end{proof}

\begin{remark}[Scope of the converse]
The quantifier ``fixed $m$'' matters. The proof uses a fixed meromorphic multiplier and does not prove a converse for a depth varying with $t$. Equal Bernoulli weights matter too: a factor $p+(1-p)e^{hw}$ with $p>1/2$ has zeros at real part $h^{-1}\log(p/(1-p))>0$ and could cancel a selected off-line pole. Our filter's zeros stay on the allowed boundary lines. Real positive residues are why an \emph{upper} bound is used; no absence of real zeros (e.g. a hypothetical exceptional zero) is assumed.
\end{remark}

\section{Consequences and variants}
\begin{corollary}[Central zeros from one starting scale]\label{cor:one}
Under $\GRH_X$,
\[
 M=\frac14\lim_{m\to\infty}B_m d(0).
\]
Each expression uses only prime powers $n\le\exp(1+(1+\sqrt2)m)$. The limit, and the eventual integer rounding, require no tabulated zero ordinates and no separately estimated constant in the prime number theorem.
\end{corollary}
\begin{proof}
Set $t=0$ in Theorem \ref{th:main}; the largest shift is $m(1+\sqrt2)$ and \eqref{eq:d} evaluates at one further unit.
\end{proof}

\begin{corollary}[Isolating individual real characters]\label{cor:real}
Let $\chi$ be a nonprincipal real Dirichlet character modulo $q$ and let $H=\ker\chi$. Under GRH for the principal and $\chi$, the sequence $B_m d(0)/4$ converges to $\ord_{1/2}L(s,\chi)$. Thus the individual central multiplicity can be recovered from prime powers with $\chi(n)=1$, with the normalization $[G:H]=2$.
\end{corollary}
\begin{proof}
$H^\perp=\{1,\chi\}$. The principal $L$-function is zeta times Euler factors nonzero at $1/2$, and $\zeta(1/2)\ne0$. The latter follows since the alternating eta series is strictly positive at $1/2$ and $1-2^{1/2}<0$. Apply Corollary \ref{cor:one}.
\end{proof}

For a complex character $\chi$ of order $k$, one can apply the construction to $H_d=\ker(\chi^{k/d})$ for each $d\mid k$. These kernels have annihilator $\langle\chi^{k/d}\rangle$, the unique order-$d$ subgroup of $\langle\chi\rangle$. If $M_d$ denotes the associated total central multiplicity, then
\[
 \sum_{\substack{1\le j\le k\\(j,k)=1}}\ord_{1/2} L(s,\chi^j)
 =\sum_{d\mid k}\mu(k/d) M_d
\]
by ordinary M\"obius inversion. This gives the sum over generators, not each complex character separately. Similarly, if $M_q$ denotes the total for all characters modulo $q$, and $P_q$ the total for primitive characters of exact conductor $q$, then $M_q=\sum_{d\mid q}P_d$ and $P_q=\sum_{d\mid q}\mu(q/d)M_d$. Induction from a primitive character changes only nonvanishing Euler factors at $1/2$.

\begin{proposition}[Exact finite-band annihilation]\label{prop:targeted}
If one is willing to use zero ordinates, the filter can instead eliminate every noncentral critical-line ordinate of absolute value at most $U\ge1$ exactly. Under $\GRH_X$, form the positive averaging operator
\[
 C_U=\prod_{0<\gamma\le U}^{\mathrm{distinct}} A_{\pi/\gamma},
\]
where the product uses positive ordinates of the family. Then $C_U d(t)=4M+O_{q,H}(\log(2U)/U)+O_{q,H}(e^{-t/4})$, uniformly in $t\ge0$. Its maximal translation is $O_{q,H}(\log^2(2U))$. The analogue of Theorem \ref{th:converse} holds for every such finite product (and indeed for any fixed finite product of positive-step Bernoulli operators).
\end{proposition}
\begin{proof}
The family is conjugation-stable, so the positive ordinates cover the nonzero pairs. Each factor annihilates its designated pair $\pm\gamma$ by $1+e^{\pm i\pi}=0$. The remaining multipliers on the imaginary axis have modulus at most one, and the zero-counting tail is $O_{q,H}(\log(2U)/U)$ by dyadic summation. The exponentially decaying error retains its envelope because shifts are nonnegative. Summing $1/\gamma$ over the positive ordinates up to $U$ gives $O_{q,H}(\log^2(2U))$ by dyadic summation and the positive gap from zero. For the converse the new multiplier has no zeros in $0<\Rea w<1/2$ and is positive on the corresponding real interval; the proof of Theorem \ref{th:converse} applies verbatim.
\end{proof}

\begin{remark}[A limit of the two-step argument]
The inverse-square exponent in Lemma \ref{lem:cos} cannot be improved uniformly over arbitrary real frequencies. There are arbitrarily large positive integers $a,b$ with $b^2-2a^2=\pm1$, from powers of $1+\sqrt2$. At $u=2\pi a$, the first cosine has modulus one and $|\sqrt2a-b|=(\sqrt2a+b)^{-1}=O(1/a)$, so the second cosine has modulus $1-O(a^{-2})$. A uniform damping bound with exponent $c/(1+|u|)^p$ for $p<2$ is therefore impossible. This does not show that the aggregate bound \eqref{eq:main} is sharp for actual zero ordinates.
\end{remark}

\section{Numerical illustration and verification status}
The accompanying \texttt{experiment.py} evaluates $S(x)$ exactly as a finite prime-power sum, with floating-point logarithms and arithmetic. It enumerates prime powers by a sieve, constructs prefix sums of $a(n)$ and $a(n)/n$, then evaluates $S(x)$ as their difference after division by $x$. For all test moduli it cross-checks against an independent trial-division implementation at several small arguments. The run generated 60 rows (four moduli, three base scales, five depths), using 834,329 prime powers up to 12,744,840. Table \ref{tab:results} shows the identity-subgroup case at $x=e^t=300$. These numbers are illustrative, not interval-certified or evidence of a proof of GRH; the decrease need not be monotone in other examples. No assertion of central nonvanishing for a nonprincipal character is inferred from these small-depth tests.

\begin{table}[ht]
\centering
\caption{Computed $B_m d(\log 300)$ for the full character group modulo $q$ (rounded).}\label{tab:results}
\begin{tabular}{rrrrrr}\toprule
$q$&$m=0$&$m=1$&$m=2$&$m=3$&$m=4$\\\midrule
1&0.196567&0.120599&0.063935&0.031444&0.019743\\
3&0.793994&0.531388&0.328310&0.199614&0.126026\\
5&1.417600&0.799197&0.535496&0.395121&0.271215\\
8&0.974249&0.829372&0.488385&0.329996&0.221383\\\bottomrule
\end{tabular}
\end{table}

\paragraph{Proof status.} The deductions are classical analytic proofs using standard Dirichlet $L$-function analytic continuation, nonvanishing on the line $1$, functional equations, completed Hadamard products and zero-counting estimates, plus the stated GRH for the forward direction only. The converse is unconditional. The finite computation is a sanity check of the statistic, not a substitute for any lemma. No Lean formalization of the new results is claimed: the local environment has no Lean executable and the repository build requires its pinned toolchain and a substantial Mathlib installation. A formalization roadmap and dependency ledger are included separately.

\section{Further research questions}
\begin{enumerate}
\item Can one bound the constants effectively and uniformly in $q,H$, using explicit zero counts and certified lower bounds for the least positive ordinate? The present theorem is fixed-family and has no explicit rounding threshold.
\item Can higher-dimensional badly approximable translation vectors or nonuniform step schedules improve the error for a given number of sample points? Any proposed filter should retain nonvanishing in the open strip $0<\Rea w<1/2$ and positivity on its real segment.
\item Can one recover individual nonreal-character central orders using a finite system of real prime statistics without sacrificing the one-sided detection property for the whole family? The subgroup construction only yields generator-sums.
\item Can the fixed-depth converse be upgraded to slowly increasing depth as $t\to\infty$? The present Laplace multiplier argument does not immediately apply.
\item Can the same scheme be used for positive-coefficient families of other $L$-functions after establishing suitable zero locations? Such applications would need genuinely broader GRH inputs, not provided by this repository.
\end{enumerate}

\begin{thebibliography}{9}\small
\bibitem{repo} L. Sernau et al., \emph{GRH for Dirichlet L-functions, in Lean 4}, 25 Nov. 2026, commit \texttt{72fead1}. \url{https://github.com/lsernau/grh-lean/tree/72fead102f663d5c7a407964a13dc0cf9d977207}.
\bibitem{rubinstein} M. Rubinstein, \emph{Alp\"oge/Claude proof of the Generalized Riemann Hypothesis}, verification and explanatory notes, Nov. 2026. \url{https://www.math.uwaterloo.ca/~mrubinst/alpoge_claude_grh/}.
\bibitem{wiki} Wikipedia contributors, \emph{Riemann hypothesis}, especially the consequences and Dirichlet $L$-series sections, accessed 26 Nov. 2026. \url{https://en.wikipedia.org/wiki/Riemann_hypothesis}.
\bibitem{mv} H. L. Montgomery and R. C. Vaughan, \emph{Multiplicative Number Theory I: Classical Theory}, Cambridge University Press, 2007, especially the explicit formula and oscillation chapters (\S15.1--15.2 for Landau-type principles).
\bibitem{dav} H. Davenport, \emph{Multiplicative Number Theory}, 3rd ed., revised by H. L. Montgomery, Springer, 2000. Classical analytic theory of Dirichlet $L$-functions and their zeros.
\bibitem{bennett} M. A. Bennett, G. Martin, K. O'Bryant and A. Rechnitzer, Counting zeros of Dirichlet $L$-functions, \emph{Math. Comp.} \textbf{90} (2021), 1455--1482. \url{https://doi.org/10.1090/mcom/3599}.
\bibitem{suzuki} M. Suzuki, On variants of Chebyshev's conjecture, \emph{Ramanujan J.} \textbf{68} (2025), article 95. \url{https://doi.org/10.1007/s11139-025-01238-9}; \url{https://arxiv.org/abs/2411.07436}.
\bibitem{pr} D. J. Platt and O. Ramar\'e, \emph{Explicit estimates: from $\Lambda(n)$ in arithmetic progressions to $\Lambda(n)/n$}, preprint version dated 5 March 2015. \url{https://research-information.bris.ac.uk/ws/files/67056122/PsiTilde_12_2.pdf}.
\end{thebibliography}
\end{document}
